% section: 議論のまとめ

\section{議論のまとめ}\label{節：議論のまとめ}
  ここまでの議論が長くなったので,扱った内容をまとめる.

  複素数$z=x+iy$を座標平面上の点$(x,y)$に対応させると,複素数を平面上の点として考えられる.
  極形式
  \[
    z=\rho(\cos\theta+i\sin\theta)
  \]
  では,\,$\rho$が原点からの距離を表し,\,$\theta$が偏角を表す.

  複素数に一定の数を加えることは平行移動に対応し,極形式で表した複素数を掛けることは回転と拡大縮小に対応する.
  積を繰り返すと偏角が加わるため,ド・モアブルの定理が得られる.
  また,第\ref{節：極限}節で準備した極限を用いるとオイラーの公式
  \[
    e^{i\theta}=\cos\theta+i\sin\theta
  \]
  が導かれる.
  このように,複素数を極形式で表すことで,座標平面上の変換と複素数の計算を結び付けられる.
